\documentclass[11pt,a4paper]{article}

\usepackage[utf8]{inputenc}
\usepackage[T1]{fontenc}
\usepackage{mathptmx}
\usepackage[margin=1in]{geometry}
\usepackage{graphicx}
\usepackage{booktabs}
\usepackage{array}
\usepackage{hyperref}
\usepackage{xcolor}
\usepackage{float}
\usepackage{amsmath,amssymb}
\usepackage{natbib}
\usepackage{fancybox}
\usepackage{setspace}
\usepackage{caption}
\usepackage{multirow}
\usepackage{longtable}
\usepackage{tabularx}

\hypersetup{
    colorlinks=true,
    linkcolor=blue,
    citecolor=blue,
    urlcolor=blue,
    pdftitle={A Conservation Constraint in LLM Context Processing Across Four Epistemological Domains},
    pdfauthor={Laxman M M},
}

% ================================================================
\title{\LARGE\textbf{A Conservation Constraint in LLM Context Processing\\
Across Four Epistemological Domains}}

\author{
    \textbf{Dr.\ Laxman M M, MBBS}\\
    Government Duty Medical Officer, PHC Manchi\\
    Bantwal Taluk, Dakshina Kannada, Karnataka, India\\
    DNB General Medicine Resident (2026), KC General Hospital, Bangalore\\
    Email: \texttt{barlax5377@gmail.com}\\
    ORCID: \href{https://orcid.org/0009-0009-0405-6531}{0009-0009-0405-6531}
}

\date{April 2026}

% ================================================================
\begin{document}

\maketitle

\begin{center}
\fbox{\parbox{0.9\textwidth}{
\textbf{Paper 6 of the MCH Research Program} --- The capstone of a nine-paper programme
measuring context sensitivity and output variance in large language models across domains.
We report an empirical conservation constraint: the product of context sensitivity
($\Delta$RCI) and output variance (Var\_Ratio) is approximately constant within an
epistemological domain across 24 model-domain runs, 14 architectures, 8 vendors, and
three independent embedding models.
}}
\end{center}

\vspace{0.5em}

% ================================================================
\begin{abstract}
\noindent We report an empirical conservation constraint in large language model context
processing: the product of context sensitivity ($\Delta$RCI) and output variance
(Var\_Ratio) is approximately constant within an epistemological domain. This constraint
--- $K = \Delta\text{RCI} \times \text{Var\_Ratio}$ --- holds across 24 model-domain
runs spanning 14 architectures from 8 vendors, four knowledge domains representing four
truth-types (Medical/Discovered, Legal/Argued, Philosophy/Explored, Applied Ethics/Felt),
and three independent embedding models (384-dimensional MiniLM, 768-dimensional MPNet,
768-dimensional LaBSE). Under the primary embedding model (all-MiniLM-L6-v2),
$K(\text{Medical}) = 0.429$ (CV $= 0.170$, $N = 8$), $K(\text{Legal}) = 0.348$
(CV $= 0.214$, $N = 5$), $K(\text{Philosophy}) = 0.301$ (CV $= 0.166$, $N = 6$),
$K(\text{Ethics}) = 0.223$ (CV $= 0.162$, $N = 5$). The conservation principle persists
under two additional embedding models, though absolute $K$ values are
embedding-scale-dependent and domain ordering shows partial sensitivity to measurement
choice. Under the primary embedding, Medical (discovered truth) yields the highest $K$;
Ethics (felt truth) yields the lowest. Medical remains consistently highest or near-highest
across all three embeddings; Ethics consistently among the lowest. Each truth-type determines a distinct behavioural
mode characterised by entanglement structure, exploration arc, content-order balance, and
variance signature. The conservation constraint is more fundamental than entanglement,
which is domain- and embedding-dependent, and more robust than temporal dynamics, which
vary by model. The constraint holds across three independent embedding models; whether it holds
under non-English encoding degradation (EFI $= 0.07$ for Kannada, Paper 8) is
unknown as $\Delta$RCI was not measured for non-English inputs. These findings establish that domain epistemology, not
model architecture, is the primary determinant of context processing structure in LLMs.

\vspace{0.5em}
\noindent\textbf{Keywords:} large language models, context sensitivity, conservation
constraint, epistemological domains, behavioural modes, output variance, $\Delta$RCI,
embedding robustness

\vspace{0.5em}
\noindent\textbf{OSF Pre-registration:} \url{https://osf.io/dp8nj/}\\
\noindent\textbf{Data and Code:} \url{https://github.com/LaxmanNandi/MCH-Research}
\end{abstract}

% ================================================================
\section{Introduction}

\subsection{Context Sensitivity in Large Language Models}

Context sensitivity --- how a model's output changes based on conversational history ---
is a fundamental property of large language models (LLMs) affecting deployment, safety,
and reliability \citep{liang2023helm,chang2024survey}. A model that responds identically
with or without prior context provides no genuine contextual grounding; a model that is
maximally context-sensitive may be unstable or vulnerable to manipulation. Empirical work
has shown that models do not use context uniformly --- performance degrades when relevant
information is positioned in the middle of long inputs rather than at the edges
\citep{liu2024lost,bai2024longbench,zhang2024infinitebench}, and models effectively
utilise only 10--20\% of available context on complex reasoning tasks
\citep{kuratov2024,hsieh2024ruler,an2025context}. Despite extensive work on model capabilities
\citep{brown2020,kaplan2020,singhal2023,hendrycks2021,chiang2024chatbot} and alignment
\citep{hubinger2019,ouyang2022,bai2022constitutional}, and systematic evaluation
\citep{zheng2023,sclar2024sensitivity}, the \emph{structure}
of context sensitivity across epistemological domains has received limited systematic
attention.

Two observations motivate the present study. First, earlier papers in this programme
\citep{laxman2026a,laxman2026b} established that context sensitivity, measured as
$\Delta\text{RCI} = \text{RCI}_\text{TRUE} - \text{RCI}_\text{COLD}$, varies
substantially across both models and domains. Second, Paper 4 \citep{laxman2026d}
reported a striking correlation between context sensitivity and variance reduction
($r = 0.76$, $p < 0.001$) in the Medical domain --- suggesting the two quantities
are not independent. Large-scale simulation work has confirmed that LLMs exhibit
systematic performance degradation in multi-turn conversations, with unreliability
increasing by 112\% across 200,000+ conversations \citep{laban2025,kwan2024mteval,wang2024mint,maharana2024memory} --- a phenomenon
Paper 4 showed corresponds to divergent entanglement ($\text{Var\_Ratio} > 1$). The
present paper asks whether a conservation constraint governs the relationship between
$\Delta$RCI and output variance across domains.

\subsection{The Research Programme}

This paper is the capstone of a nine-paper programme. Papers 1--2
\citep{laxman2026a,laxman2026b} established the $\Delta$RCI metric and a standardised
experimental protocol across 14 model architectures. Paper 3 \citep{laxman2026c}
identified domain-specific temporal dynamics, including the U-shape (Medical) and
inverted-U (Philosophy) patterns across 30-prompt conversations. Paper 4
\citep{laxman2026d} discovered entanglement between context sensitivity and variance
reduction ($r = 0.76$) in the Medical domain. Paper 5 \citep{laxman2026e} built a
safety taxonomy (IDEAL/EMPTY/DIVERGENT/RICH) based on $\Delta$RCI and Var\_Ratio
profiles. Paper 7 \citep{laxman2026g} decomposed $\Delta$RCI into content and order
components using a SCRAMBLED condition. Paper 8 \citep{laxman2026h} measured encoding
fidelity failure for non-English languages (EFI $= 0.07$ for Kannada), finding that
the conservation constraint holds even under near-complete semantic degradation. Paper 9
\citep{laxman2026i} validated that encoding fidelity assessment is embedding-model
dependent, establishing LaBSE as the most semantically faithful measurement instrument.

In the course of this programme, we observed that the product $\Delta\text{RCI} \times
\text{Var\_Ratio}$ clusters within domains despite varying substantially across
architectures. This paper formalises and validates this observation across four
epistemological domains, 24 model-domain runs, and three independent embedding models.

\subsection{Contributions}

This paper makes five contributions:

\begin{enumerate}
    \item We report $K = \Delta\text{RCI} \times \text{Var\_Ratio} \approx$ constant
    within a domain, validated across 24 model-domain runs spanning 14 architectures
    from 8 vendors.
    \item We identify four behavioural modes corresponding to four truth-types
    (Discovered, Argued, Explored, Felt), each with a distinct signature of entanglement,
    arc, content fraction, and variance structure.
    \item We demonstrate that the conservation principle holds across three independent
    embedding models, though absolute $K$ values are embedding-scale-dependent.
    \item We show that conservation ($K$) is more fundamental than entanglement, which
    is domain- and embedding-dependent, and that domain epistemology is the primary
    determinant of context processing mode.
    \item We establish that the Var\_Ratio component of $K$ is elevated under
    non-English encoding degradation (Kannada VR $= 2.36$--$4.08$, Paper 8
    \citep{laxman2026h}), though direct measurement of $\Delta$RCI in non-English
    conditions remains a target for future work.
\end{enumerate}

% ================================================================
\section{Methods}

\subsection{Experimental Protocol}

All experiments followed the standardised MCH protocol established in Papers 1--2
\citep{laxman2026a,laxman2026b}. Three conditions were administered per model-domain
configuration:

\begin{itemize}
    \item \textbf{TRUE}: Coherent 29-message conversational history presented before
    each prompt. Context is semantically and temporally coherent.
    \item \textbf{COLD}: No context --- each prompt presented as a standalone query.
    \item \textbf{SCRAMBLED}: Randomised prompt order with a coherent conversation
    header. Content is present; order is destroyed.
\end{itemize}

Each condition comprised 50 trials per model-domain configuration, with 30 prompts per
trial. Temperature was set to 0.7 for all models to permit variance measurement
\citep{renze2024,shi2024decoding,atil2024,dietrich2025reproducibility}. Total responses per domain: Medical
$8 \times 50 \times 30 \times 3 = 36{,}000$; aggregated programme total exceeding
160,000 responses.

\subsection{Metrics}

\begin{itemize}
    \item $\Delta\text{RCI} = \overline{\text{RCI}_\text{TRUE}} -
    \overline{\text{RCI}_\text{COLD}}$: context sensitivity, where RCI is the
    mean pairwise cosine similarity across trials at each prompt position.
    \item $\text{Var\_Ratio} = \text{Var}(\text{TRUE cross-trial pairwise similarity})
    / \text{Var}(\text{COLD cross-trial pairwise similarity})$: output variance ratio
    under context versus without context.
    \item $K = \Delta\text{RCI} \times \text{Var\_Ratio}$: the conservation product.
    \item $\text{VRI} = 1 - \text{Var\_Ratio}$: Variance Reduction Index, used in
    entanglement analysis.
    \item Content Fraction $= \overline{\text{SCRAM}-\text{COLD}} /
    \overline{\text{TRUE}-\text{COLD}} \times 100\%$: percentage of $\Delta$RCI
    attributable to content alone (from Paper 7 \citep{laxman2026g}).
    \item Exploration Arc $= \text{Var}(\text{TRUE at P30}) /
    \text{Var}(\text{TRUE at P1})$: convergence ($<1$) or divergence ($>1$) of
    response variance over the conversation.
    \item Entanglement: Pearson $r$ between position-level $\Delta$RCI and VRI
    across 30 prompt positions within a model-domain run.
\end{itemize}

\subsection{Domains and Truth-Types}

Four epistemological domains were selected to represent distinct types of knowledge
validity (Table~\ref{tab:domains}).

\begin{table}[H]
\centering
\caption{Domains, truth-types, and prompt descriptions.}
\label{tab:domains}
\begin{tabular}{llll}
\toprule
\textbf{Domain} & \textbf{Truth-Type} & \textbf{Prompts} & \textbf{Structure} \\
\midrule
Medical     & Discovered & 30 (STEMI case)          & Fixed diagnostic answer exists \\
Legal       & Argued     & 30 (whistleblower case)  & Answer constructed through reasoning \\
Philosophy  & Explored   & 30 (consciousness)       & Open inquiry, no fixed answer \\
Appl.\ Ethics & Felt     & 30 (trolley to synthesis) & Moral commitment; context individuates \\
\bottomrule
\end{tabular}
\end{table}

The truth-type classification reflects the epistemological basis of validity in each
domain: Medical knowledge converges toward discoverable ground truth
\citep{thirunavukarasu2023,singhal2025expert,ayers2023physician,nori2023}; Legal knowledge
is constructed through argumentation within a rule system
\citep{guha2023,nay2024tax,dahl2024hallucinations,choi2022lawschool}; Philosophical
knowledge expands through inquiry without requiring closure; Ethical knowledge is grounded
in values that resist purely logical resolution
\citep{shanahan2024,fierro2024,bender2021,ji2024moralbench,suzgun2024belief,paseri2025epistemological}.

\subsection{Models}

Twenty-four model-domain runs were completed across 14 unique architectures from 8
vendors (Table~\ref{tab:models_full}). Models were accessed via API: Together.ai for
open-weight models, vendor APIs for closed models. The model set spans dense transformers
\citep{vaswani2017} and mixture-of-experts architectures \citep{fedus2022switch},
representing parameter counts from 14B to 1T \citep{hoffmann2022,wei2022emergent,kaplan2020}.

\begin{table}[H]
\centering
\caption{All 24 model-domain runs with architecture details.}
\label{tab:models_full}
\begin{tabular}{llll}
\toprule
\textbf{Model} & \textbf{Parameters} & \textbf{Vendor} & \textbf{Domains} \\
\midrule
DeepSeek V3.1         & 671B (37B active)  & DeepSeek     & Medical, Philosophy, Legal, Ethics \\
Qwen3 235B            & 235B (22B active)  & Alibaba      & Medical, Legal \\
Llama 4 Maverick      & 400B (17B active)  & Meta         & Medical, Philosophy, Ethics \\
Llama 4 Scout         & 109B (17B active)  & Meta         & Medical \\
Llama 3.3 70B Turbo   & 70B               & Meta         & Legal \\
Mistral Small 24B     & 24B               & Mistral      & Medical, Legal, Ethics \\
Ministral 14B         & 14B               & Mistral      & Medical \\
Kimi K2               & 1T (32B active)    & Moonshot     & Medical \\
Claude Haiku          & ---               & Anthropic    & Philosophy \\
GPT-4o                & ---               & OpenAI       & Philosophy \\
GPT-4o-mini           & ---               & OpenAI       & Philosophy \\
Gemini Flash          & ---               & Google       & Medical, Philosophy \\
\bottomrule
\end{tabular}
\end{table}

\subsection{Embedding Models}

Three embedding models were used for robustness validation (Table~\ref{tab:embeddings}):

\begin{table}[H]
\centering
\caption{Embedding models used for $K$ computation.}
\label{tab:embeddings}
\begin{tabular}{llll}
\toprule
\textbf{Model} & \textbf{Dim} & \textbf{Category} & \textbf{Reference} \\
\midrule
all-MiniLM-L6-v2  & 384 & EN-distilled      & \citet{reimers2020} \\
all-mpnet-base-v2 & 768 & EN-distilled      & \citet{reimers2020} \\
LaBSE             & 768 & Language-agnostic & \citet{feng2022} \\
\bottomrule
\end{tabular}
\end{table}

LaBSE was validated as semantically faithful by Paper 9 \citep{laxman2026i}: parallel
versus non-parallel sentence pair discrimination $\Delta = 0.58$, $p = 0.004$
(Mann-Whitney U, $N = 5$ per group). Cross-lingual representation quality varies
substantially across language families \citep{conneau2020,jin2024}. All-MiniLM-L6-v2
serves as the primary embedding for the main analysis; MPNet and LaBSE provide
robustness checks. Embedding model selection followed the MTEB benchmark
\citep{muennighoff2023,enevoldsen2025mmteb}, which established that no single embedding
method dominates across tasks, motivating multi-model validation. Cosine similarity as
a semantic proxy has been examined by \citet{gao2021simcse,steck2024cosine,ethayarajh2019contextual}.

\subsection{Re-embedding Protocol}

For Var\_Ratio computation under MPNet and LaBSE, TRUE condition alignments previously
stored as 1.0 (self-comparison artefact) were re-computed via cross-trial pairwise
cosine similarity at each prompt position. This re-embedding was performed independently
for all three embedding models using the stored LLM response text, ensuring comparability
across the three-embedding robustness analysis.

\subsection{Statistical Tests}

Domain separation was tested using Mann-Whitney $U$ (non-parametric, two-sided). Effect
sizes reported as Cohen's $d$. Coefficient of variation (CV $=$ SD/mean) used to assess
within-domain clustering tightness. Entanglement tested as Pearson $r$ across 30 prompt
positions, with significance threshold $\alpha = 0.05$.

\subsection{AI-Assisted Tools}

AI tools (Claude, Anthropic; ChatGPT, OpenAI; DeepSeek) were used to assist with data
analysis, statistical computation, visualisation, and manuscript preparation. All
experimental design, clinical interpretations, hypothesis formulation, and scientific
conclusions are the sole responsibility of the author.

% ================================================================
\section{Results}

\subsection{Conservation Constraint Under Primary Embedding}

Table~\ref{tab:all_runs} presents all 24 model-domain runs with $\Delta$RCI, Var\_Ratio,
$K$, entanglement $r$, and exploration arc under the primary embedding model
(all-MiniLM-L6-v2).

\begin{longtable}{llcccc}
\caption{All 24 model-domain runs: $\Delta$RCI, Var\_Ratio, $K$, entanglement $r$, and
arc (primary embedding: all-MiniLM-L6-v2).}
\label{tab:all_runs}\\
\toprule
\textbf{Model} & \textbf{Domain} & \textbf{$\Delta$RCI} & \textbf{Var\_Ratio} &
\textbf{$K$} & \textbf{Ent.\ $r$} \\
\midrule
\endfirsthead
\toprule
\textbf{Model} & \textbf{Domain} & \textbf{$\Delta$RCI} & \textbf{Var\_Ratio} &
\textbf{$K$} & \textbf{Ent.\ $r$} \\
\midrule
\endhead
\midrule
\multicolumn{6}{r}{\textit{Continued on next page}} \\
\endfoot
\bottomrule
\endlastfoot
% Medical
Gemini Flash         & Medical    & 0.427 & 1.287 & 0.549 & 0.389* \\
Llama 4 Scout        & Medical    & 0.323 & 1.610 & 0.521 & 0.473** \\
Qwen3 235B           & Medical    & 0.328 & 1.334 & 0.437 & 0.655*** \\
Ministral 14B        & Medical    & 0.391 & 1.080 & 0.423 & 0.649*** \\
Kimi K2              & Medical    & 0.417 & 1.006 & 0.420 & 0.653*** \\
Llama 4 Maverick     & Medical    & 0.316 & 1.213 & 0.384 & 0.624*** \\
Mistral Small 24B    & Medical    & 0.365 & 0.985 & 0.359 & 0.780*** \\
DeepSeek V3.1        & Medical    & 0.320 & 1.071 & 0.343 & 0.672*** \\
\midrule
% Philosophy
Gemini Flash         & Philosophy & 0.338 & 1.120 & 0.378 & 0.572*** \\
Claude Haiku         & Philosophy & 0.331 & 1.012 & 0.334 & 0.584*** \\
DeepSeek V3.1        & Philosophy & 0.302 & 1.034 & 0.312 & 0.447* \\
GPT-4o               & Philosophy & 0.283 & 0.950 & 0.269 & 0.564** \\
GPT-4o-mini          & Philosophy & 0.269 & 0.968 & 0.260 & 0.582*** \\
Llama 4 Maverick     & Philosophy & 0.266 & 0.939 & 0.250 & 0.644*** \\
\midrule
% Legal
Qwen3 235B           & Legal      & 0.265 & 1.796 & 0.476 & ns \\
DeepSeek V3.1        & Legal      & 0.276 & 1.225 & 0.338 & ns \\
Mistral Small 24B    & Legal      & 0.252 & 1.338 & 0.338 & ns \\
Llama 4 Maverick     & Legal      & 0.209 & 1.412 & 0.295 & ns \\
Llama 3.3 70B Turbo  & Legal      & 0.206 & 1.428 & 0.294 & ns \\
\midrule
% Ethics
DeepSeek V3.1        & Ethics     & 0.257 & 1.058 & 0.272 & ns \\
Qwen3 235B           & Ethics     & 0.211 & 1.184 & 0.250 & ns \\
Llama 4 Maverick     & Ethics     & 0.181 & 1.176 & 0.213 & ns \\
Llama 3.3 70B Turbo  & Ethics     & 0.172 & 1.128 & 0.194 & ns \\
Mistral Small 24B    & Ethics     & 0.183 & 1.029 & 0.188 & ns \\
\end{longtable}

\noindent Significance codes: * $p < 0.05$, ** $p < 0.01$, *** $p < 0.001$.

Domain-level summary statistics are presented in Table~\ref{tab:domain_summary}.

\begin{table}[H]
\centering
\caption{Conservation constraint by domain (primary embedding: all-MiniLM-L6-v2).}
\label{tab:domain_summary}
\begin{tabular}{lcccccc}
\toprule
\textbf{Domain} & \textbf{N} & \textbf{K mean} & \textbf{K SD} & \textbf{CV} &
\textbf{K min} & \textbf{K max} \\
\midrule
Medical     & 8 & 0.429 & 0.073 & 0.170 & 0.343 & 0.549 \\
Legal       & 5 & 0.348 & 0.075 & 0.214 & 0.294 & 0.476 \\
Philosophy  & 6 & 0.301 & 0.050 & 0.166 & 0.250 & 0.378 \\
Ethics      & 5 & 0.223 & 0.036 & 0.163 & 0.188 & 0.272 \\
\bottomrule
\end{tabular}
\end{table}

Domain separation is statistically significant across five of six pairwise comparisons
(Table~\ref{tab:pairwise}). The domain ordering under MiniLM is: Discovered $>$ Argued
$>$ Explored $>$ Felt. CVs range from 0.162 (Ethics) to 0.214 (Legal), indicating
consistent within-domain clustering despite large variation across architectures
(14B to 1T parameters). Philosophy vs Legal is not significant ($p = 0.247$), consistent
with their adjacent $K$ values (0.301 vs 0.348).

\begin{table}[H]
\centering
\caption{Pairwise Mann-Whitney $U$ tests between domains (primary embedding: MiniLM).}
\label{tab:pairwise}
\begin{tabular}{llll}
\toprule
\textbf{Comparison} & \textbf{p-value} & \textbf{Cohen's d} & \textbf{Significant} \\
\midrule
Medical vs Ethics      & 0.002 & 3.57 & Yes *** \\
Medical vs Philosophy  & 0.003 & 2.06 & Yes ** \\
Legal vs Ethics        & 0.008 & 2.13 & Yes ** \\
Philosophy vs Ethics   & 0.030 & 1.77 & Yes * \\
Medical vs Legal       & 0.045 & 1.10 & Yes * \\
Philosophy vs Legal    & 0.247 & ---  & No (ns) \\
\bottomrule
\end{tabular}
\end{table}

\begin{figure}[H]
\centering
\includegraphics[width=0.95\textwidth]{fig1_conservation_scatter.png}
\caption{Conservation constraint scatter plot: all 24 model-domain runs. X-axis:
$\Delta$RCI. Y-axis: Var\_Ratio. Each domain colour-coded (Medical: blue, Legal: green,
Philosophy: orange, Ethics: red). Domain-specific hyperbolas ($K = $ constant curves)
overlaid. Models cluster along their respective domain hyperbolas despite spanning
14 architectures from 8 vendors.}
\label{fig:conservation_scatter}
\end{figure}

\begin{figure}[H]
\centering
\includegraphics[width=0.85\textwidth]{fig2_k_by_domain.png}
\caption{$K$ values by domain. Box plots showing means, interquartile ranges, and
individual model values. CV shown above each domain. Domains are clearly separated,
with Medical consistently highest (mean $K = 0.429$) and Ethics consistently lowest
(mean $K = 0.223$). Domain ordering: Discovered $>$ Argued $>$ Explored $>$ Felt
under primary embedding.}
\label{fig:k_by_domain}
\end{figure}

\subsection{Four Behavioural Modes}

Each truth-type determines a distinct behavioural signature across four dimensions:
entanglement structure, exploration arc, content fraction, and variance profile
(Table~\ref{tab:modes}).

\begin{table}[H]
\centering
\caption{Behavioural mode signatures across four truth-types.}
\label{tab:modes}
\begin{tabular}{lllllll}
\toprule
\textbf{Mode} & \textbf{K} & \textbf{Ent.} & \textbf{Arc} & \textbf{Content\%} &
\textbf{VR} & \textbf{P30 Spike} \\
\midrule
Discovered (Med.)  & 0.429 & Yes ($r=0.76$) & Conv.\ (1.72) & 45--55\% & $>1$ & Yes \\
Argued (Legal)     & 0.348 & No (MiniLM)   & Conv.         & 70--80\% & $\approx 1$ & No \\
Explored (Phil.)   & 0.301 & Yes ($r=0.76$) & Div.\ (15.23) & 35--55\% & $\approx 1$ & No \\
Felt (Ethics)      & 0.223 & Mixed         & Mixed         & 80--91\% & $<1$ to $\approx1$ & No \\
\bottomrule
\end{tabular}
\end{table}

\begin{figure}[H]
\centering
\includegraphics[width=0.85\textwidth]{fig3_four_mode_taxonomy.png}
\caption{Four-mode taxonomy: 2$\times$2 quadrant plot. X-axis: Entangled vs Not
Entangled (under MiniLM). Y-axis: Convergent vs Divergent arc. Medical occupies the
entangled-convergent quadrant. Philosophy occupies the entangled-divergent quadrant.
Legal occupies the not-entangled-convergent quadrant. Ethics occupies the mixed zone.
Note: Legal entanglement status is embedding-dependent (see Section~\ref{sec:entanglement}).}
\label{fig:four_mode}
\end{figure}

\subsection{Robustness: Three Embedding Models}
\label{sec:robustness}

Table~\ref{tab:three_embeddings} presents $K$ means across all three embedding models.

\begin{table}[H]
\centering
\caption{$K$ means across three independent embedding models. Conservation principle
holds across all three embeddings; absolute $K$ values and domain ordering show
partial sensitivity to embedding choice.}
\label{tab:three_embeddings}
\begin{tabular}{lcccc}
\toprule
\textbf{Domain} & \textbf{K (MiniLM 384D)} & \textbf{K (MPNet 768D)} &
\textbf{K (LaBSE 768D)} & \textbf{Consistent rank} \\
\midrule
Medical     & 0.429 & 0.188 & 0.215 & Highest (2/3) \\
Legal       & 0.348 & 0.280 & 0.193 & Variable \\
Philosophy  & 0.301 & 0.153 & 0.195 & Variable \\
Ethics      & 0.223 & 0.214 & 0.187 & Lowest (2/3) \\
\bottomrule
\end{tabular}
\end{table}

The conservation principle holds across all three embeddings in the sense that $K$
clusters within domains (CV $= 0.13$--$0.19$ under MiniLM; CV $= 0.30$--$0.67$ under
LaBSE). Absolute $K$ values are embedding-scale-dependent --- MPNet yields the smallest
$K$ values overall; MiniLM the largest. Medical is the highest or near-highest domain
across all three embeddings. Ethics is the lowest across two of three embeddings.
The middle domains (Legal and Philosophy) show ordering sensitivity to embedding choice.

\begin{figure}[H]
\centering
\includegraphics[width=0.95\textwidth]{fig4_three_embedding.png}
\caption{Three-embedding robustness: grouped bar chart. Four domains $\times$ three
embedding models. $K$ means shown for each combination. The conservation principle
(domain clustering) holds across all three embeddings, while absolute $K$ values shift
with embedding scale. Medical consistently yields the highest $K$ under MiniLM and
LaBSE.}
\label{fig:three_embeddings}
\end{figure}

\subsection{Entanglement Is Domain- and Embedding-Dependent}
\label{sec:entanglement}

Under the primary embedding (MiniLM), entanglement between $\Delta$RCI and VRI is
strong in Medical ($r = 0.76$, $p < 0.001$) and Philosophy ($r = 0.76$, $p < 0.001$),
absent in Legal ($r = -0.033$, ns), and model-dependent in Ethics.

Under LaBSE, entanglement is significant in all domains including Legal (all $r > 0.37$,
all $p < 0.05$). This means entanglement presence or absence is not a fixed property
of the domain --- it is partly determined by the embedding space used for measurement.

The conservation constraint $K$, by contrast, holds regardless of entanglement status.
$K$ is preserved in Legal under both embeddings even though entanglement is absent under
MiniLM and present under LaBSE. This establishes $K$ as more fundamental than
entanglement: the conservation budget is set by the domain; entanglement is a secondary
pattern that may or may not appear depending on measurement conditions.

\subsection{Content-Order Decomposition}

The SCRAMBLED condition separates $\Delta$RCI into content and order components
\citep{laxman2026g}. Content fraction results are presented in
Figure~\ref{fig:content_fraction} and summarised in Table~\ref{tab:content_order}.

\begin{table}[H]
\centering
\caption{Content fraction by domain. Higher values indicate content dominates
$\Delta$RCI; lower values indicate conversational order contributes substantially.}
\label{tab:content_order}
\begin{tabular}{llll}
\toprule
\textbf{Domain} & \textbf{Content\%} & \textbf{Interpretation} & \textbf{K} \\
\midrule
Ethics      & 80--91\% & Content dominates; order IS moral reasoning & 0.223 \\
Legal       & 70--80\% & Content dominates; framework reduces order dependence & 0.348 \\
Medical     & 45--55\% & Both content and order contribute equally & 0.429 \\
Philosophy  & 35--55\% & Order contributes substantially; inquiry is sequential & 0.301 \\
\bottomrule
\end{tabular}
\end{table}

The inverse relationship between content fraction and $K$ is notable: domains where
order contributes more (Medical, Philosophy) show higher or comparable $K$ values.
Ethics, where SCRAMBLED $\approx$ COLD (order provides almost no signal), has the
lowest $K$.

\begin{figure}[H]
\centering
\includegraphics[width=0.85\textwidth]{fig5_content_order.png}
\caption{Content fraction gradient across four domains. Bar chart showing the percentage
of $\Delta$RCI attributable to content alone (SCRAMBLED condition). Ethics (80--91\%)
shows the highest content fraction; Philosophy (35--55\%) the lowest. The gradient
reflects the degree to which conversational order carries epistemological information
in each domain.}
\label{fig:content_fraction}
\end{figure}

\subsection{Exploration Arc}

Medical domain shows a convergent arc (mean arc $= 1.72 \pm 0.68$): response variance
narrows over the 30-prompt conversation as the diagnostic case develops toward a
conclusion. Philosophy shows a strongly divergent arc (mean arc $= 15.23 \pm 16.64$):
response variance expands as inquiry deepens. The distributions of Medical and Philosophy
arcs show zero overlap, confirming that truth-type determines response topology.

Legal is convergent (similar to Medical). Ethics is mixed, with arc sign depending on
the specific model.

\subsection{K and Encoding Fidelity}

Paper 8 \citep{laxman2026h} measured EFI $= 0.07$ for Kannada clinical input ---
93\% degradation in encoding fidelity relative to English. Paper 8 also measured
variance amplification for non-English inputs: Kannada Var\_Ratio $= 2.36$--$4.08$
under LaBSE, confirmed across five embedding models \citep{laxman2026i}. This
establishes that the Var\_Ratio component of $K$ is measurably elevated for non-English
clinical input even under severe encoding degradation.

However, $\Delta$RCI --- the context sensitivity component of $K$ --- was not computed
for non-English inputs in Paper 8. The TRUE vs COLD experimental conditions were not
run in Kannada, Tamil, or Hindi. Consequently, whether the full conservation constraint
$K = \Delta\text{RCI} \times \text{Var\_Ratio}$ holds for non-English inputs is
currently unknown. The elevated Var\_Ratio under encoding degradation is a partial
signal, not a direct test of $K$.

Direct measurement of $\Delta$RCI and $K$ in non-English conditions is a target for
future work (see Section~\ref{sec:future}).

% ================================================================
\section{Discussion}

\subsection{The Conservation Constraint}

The product $\Delta\text{RCI} \times \text{Var\_Ratio}$ clusters within domains despite
enormous variation in model architecture (14B to 1T parameters), training data, vendor,
and embedding model (three independent instruments). This clustering --- summarised by
the approximately constant $K$ --- suggests a fundamental constraint on LLM context
processing: increasing context sensitivity is accompanied by decreasing output variance,
and vice versa, such that their product is domain-bounded.

The mathematical form resembles conservation laws in which the product of two varying
quantities is approximately constant. Shannon's information theory \citep{shannon1948}
established that communication systems have fundamental capacity constraints; the
information bottleneck framework
\citep{tishby2015,saxe2019,kawaguchi2023bottleneck} provides a related
precedent: it formalises a trade-off between compressing input representations and
preserving task-relevant information, structurally similar to the trade-off we observe
between context sensitivity and output variance. We do not claim the analogy is formal ---
it is structural. The relationship $K = \Delta\text{RCI} \times \text{Var\_Ratio}$
is empirically observed, not theoretically derived \citep{bahri2024scaling}. The
underlying mechanism connecting context sensitivity to output variance in transformer
architectures \citep{vaswani2017,hoffmann2022,fedus2022switch} remains unknown and is
not addressed by this paper.

\subsection{Why K Varies by Domain}

The $K$ gradient (Medical highest, Ethics lowest under MiniLM) may reflect the degree
of constraint each domain places on context processing. Discovered truth (Medical)
constrains maximally: there is a correct answer, and context drives convergence toward
it. Felt truth (Ethics) constrains minimally: each trial produces a unique moral
trajectory shaped by the specific context. The conservation budget appears to reflect
the epistemological demand imposed by the truth-type.

This interpretation is consistent with the content-order decomposition: Medical, with
balanced content and order contributions, has the highest $K$. Ethics, where order
contributes almost nothing (content fraction 80--91\%), has the lowest. Domains where
conversational order carries epistemological weight (Medical, Philosophy) show distinct
$K$ values from those where it does not (Ethics, to a lesser degree Legal).

\subsection{Robustness and Limitations of the Three-Embedding Validation}

The conservation principle holds across three independent embedding models. This
confirms that $K$ clustering is not an artifact of any single embedding space. However,
absolute $K$ values are embedding-scale-dependent (Medical $K$ ranges from 0.188 under
MPNet to 0.429 under MiniLM for the same data), and domain ordering is partially
sensitive to embedding choice --- Legal and Philosophy show nearly identical $K$ under
LaBSE ($0.193$ vs $0.195$), while they are clearly separated under MiniLM
($0.348$ vs $0.301$).

This means $K$ is best understood as a domain-specific conservation principle whose
numerical value depends on the measurement scale --- analogous to physical constants
whose numerical values depend on the unit system, while the underlying relationships
they express are scale-invariant. The within-domain CV under LaBSE
(0.30--0.67) is substantially wider than under MiniLM (0.13--0.19), indicating the
constraint is approximately rather than tightly conserved under the semantically
faithful embedding. Future work establishing $K$ on a theoretically grounded absolute
scale would strengthen the robustness claims.

\subsection{Entanglement as a Special Case}

Paper 4 \citep{laxman2026d} reported entanglement ($r = 0.76$) as a major finding
in the Medical domain. The mechanistic basis for in-context learning in transformers ---
including induction heads \citep{olsson2022induction} and chain-of-thought reasoning
\citep{wei2022cot} --- suggests that context coupling is an emergent property of
attention mechanisms \citep{wei2022emergent}. The present paper shows entanglement is
domain-dependent (absent in Legal under MiniLM) and embedding-dependent (present in
Legal under LaBSE). Conservation $K$ holds regardless of entanglement status.

This establishes a theoretical hierarchy: $K$ is more fundamental than entanglement.
The conservation budget is set by the domain epistemology. Entanglement is a secondary
pattern that appears when domain structure and embedding space jointly favour coupling
between sensitivity and variance. The hierarchy is depicted in
Figure~\ref{fig:hierarchy}.

\begin{figure}[H]
\centering
\includegraphics[width=0.75\textwidth]{fig6_hierarchy.png}
\caption{Theoretical hierarchy of constraints. Outermost level: Conservation ($K$) ---
holds across all four domains and all three embedding models. Middle level: Entanglement
--- present in Medical and Philosophy, absent in Legal under MiniLM, present in all
domains under LaBSE. Innermost level: Temporal dynamics (U-shape, inverted-U, P30 spike)
--- most model-dependent. Conservation is the most fundamental and most robust
constraint.}
\label{fig:hierarchy}
\end{figure}

\subsection{Implications for AI Safety and Deployment}

The conservation constraint provides a domain-specific expected value for $K$. Models
whose $K$ deviates substantially from the domain mean may exhibit anomalous context
processing --- either failing to use context effectively ($\Delta$RCI too low) or
producing pathologically variable outputs (Var\_Ratio too high). Paper 5's safety
taxonomy \citep{laxman2026e} (IDEAL/EMPTY/DIVERGENT/RICH) can be reframed in terms of
$K$: models in the RICH class (high $\Delta$RCI, low Var\_Ratio) occupy the expected
$K$ range; DIVERGENT models (high $\Delta$RCI, high Var\_Ratio) exceed it; EMPTY models
(low $\Delta$RCI, any Var\_Ratio) fall below it.

For clinical AI deployment specifically \citep{singhal2023,thirunavukarasu2023,pfohl2024},
the Medical domain $K = 0.429$ (CV $= 0.170$) provides a calibration benchmark. Models
showing Medical $K$ substantially outside the observed range warrant additional scrutiny
before deployment in diagnostic contexts. Output variance is a distinct safety dimension
from hallucination \citep{ji2023,lin2022,huang2024,manakul2023selfcheck} and from
sycophancy \citep{sharma2024sycophancy}: a model may produce factually correct but
highly variable outputs, or consistent but incorrect ones. Trustworthiness assessment
requires all three dimensions \citep{wang2023decodingtrust}. The conservation constraint
characterises the variance dimension specifically.

\subsection{Limitations}

\begin{enumerate}
    \item \textbf{Domain count:} Four domains with 5--8 models each. Additional domains
    (Technical, Creative, Historical) and additional models would strengthen
    generalisability.
    \item \textbf{English-language prompts only:} All prompt batteries are in English.
    Direct $K$ measurement in non-English conditions has not been performed, though
    Paper 8 \citep{laxman2026h} provides indirect evidence that $K$ survives encoding
    degradation. Multilingual evaluation gaps are well-documented
    \citep{ahuja2023,blasi2022,lai2023chatgpt,joshi2020diversity}
    and represent a priority for future work.
    \item \textbf{Approximate constraint:} CV ranges from 0.162 to 0.214 under MiniLM
    and 0.30 to 0.67 under LaBSE. The conservation is approximate; within-domain
    variation is non-trivial.
    \item \textbf{Single prompt battery per domain:} Different prompt batteries
    operationalising the same domain might yield different $K$ values. Prompt battery
    robustness has not been tested.
    \item \textbf{No theoretical derivation:} $K$ is empirically observed. The
    connection to transformer attention mechanisms \citep{vaswani2017} is unexplored.
    \item \textbf{Embedding sensitivity:} The LaBSE CVs (0.30--0.67) indicate the
    constraint weakens under the most semantically faithful embedding. Cosine similarity
    in embedding space has known geometric limitations \citep{steck2024cosine,ethayarajh2019contextual}.
    This is an honest limitation of the robustness claim.
\end{enumerate}

\subsection{Future Directions}
\label{sec:future}

\begin{enumerate}
    \item \textbf{Theoretical derivation:} What properties of transformer attention
    \citep{vaswani2017} would predict a conservation constraint of this form? A Noether-like
    analysis --- identifying the symmetry whose conservation yields $K$ --- is an
    open theoretical question.
    \item \textbf{Non-English $K$ measurement:} Direct measurement of $\Delta$RCI and
    Var\_Ratio for Kannada, Tamil, and Hindi prompts would test whether $K$ holds
    under real encoding degradation rather than the indirect evidence from Paper 8.
    Indic language resources \citep{gala2023indictrans2} and multilingual benchmarks
    \citep{ahuja2023,lai2023chatgpt} provide the infrastructure for such work.
    \item \textbf{Additional domains:} Technical (solved problems), Creative (no fixed
    answer), and Historical (discovered but past) domains would test the boundary
    conditions of the four-mode taxonomy. Legal AI benchmarks \citep{guha2023} and
    moral evaluation frameworks \citep{ji2024moralbench} are directly applicable.
    \item \textbf{LaBSE-based reanalysis:} $K$ was established primarily under MiniLM.
    Extending the full programme analysis to LaBSE \citep{enevoldsen2025mmteb} would
    establish whether the tighter constraint (CV $<$ 0.20) holds under the most
    semantically faithful embedding.
    \item \textbf{Mechanism:} The connection between $K$ and internal model representations
    --- attention patterns \citep{olsson2022induction}, layer activations, token probability
    distributions --- is entirely unexplored. If the empirical regularity reported here
    reflects a genuine structural property \citep{wei2022emergent}, it should have
    correlates in internal model states.
\end{enumerate}

The empirical regularity reported here may point toward a deeper structural principle
in how transformer architectures process contextual information. Establishing such a
principle from first principles is a target for future work; this paper establishes
the empirical basis from which such work could proceed.

% ================================================================
\section{Conclusion}

We report a conservation constraint --- $K = \Delta\text{RCI} \times \text{Var\_Ratio}$
--- that holds across four epistemological domains, 24 model-domain runs, 14
architectures from 8 vendors, and three independent embedding models. The constraint
is empirical: the product of context sensitivity and output variance is approximately
constant within a domain, while varying between domains in a pattern consistent with
the degree of epistemological constraint each truth-type imposes. Medical (discovered
truth) requires the highest conservation budget ($K = 0.429$, CV $= 0.170$); Ethics
(felt truth) the lowest ($K = 0.223$, CV $= 0.162$).

The conservation principle is more fundamental than entanglement, which is domain- and
embedding-dependent, and more robust than temporal dynamics, which vary by model and
architecture. Absolute $K$ values are embedding-scale-dependent; domain ordering shows
partial sensitivity to measurement choice. These are honest limitations that constrain
the strength of the robustness claim.

These findings establish that domain epistemology, not model architecture or parameter
count, is the primary determinant of context processing structure in LLMs. The
constraint provides a domain-specific calibration benchmark for deployment and safety
assessment. The mechanism underlying the constraint is unknown and is a target for
future theoretical work.

% ================================================================
\section*{Ethics Statement}

This study did not involve human participants, patient data, or electronic health
records. All experiments were conducted using publicly available LLM APIs with
purpose-built prompt batteries. No identifiable patient information was collected
or used at any stage. Institutional review board (IRB) approval and informed consent
were not required.

% ================================================================
\section*{Author Contributions}

L.M.M. conceived the study, designed all experiments, conducted all model-domain runs,
performed all analyses, and wrote the manuscript in its entirety. AI systems (Claude,
ChatGPT, DeepSeek) assisted with data analysis, visualisation, and manuscript
preparation under the sole direction and responsibility of the author.

% ================================================================
\section*{Competing Interests}

The author declares no competing interests.

% ================================================================
\section*{Data Availability}

All experimental data, analysis scripts, and raw results are available at
\url{https://github.com/LaxmanNandi/MCH-Research}. OSF pre-registration:
\url{https://osf.io/dp8nj/} (March 2026).

% ================================================================
\section*{Code Availability}

All analysis scripts are available at \url{https://github.com/LaxmanNandi/MCH-Research}.
Key scripts: \texttt{paper6\_conservation.py} (K computation across domains),
\texttt{paper6\_robustness\_reembed.py} (three-embedding robustness validation),
\texttt{paper6\_entanglement.py} (entanglement analysis).

% ================================================================
\section*{Acknowledgements}

No external funding was received. This work emerged from clinical practice at PHC Manchi,
Bantwal Taluk, Dakshina Kannada, Karnataka, India. AI systems (Claude, Anthropic;
ChatGPT, OpenAI; DeepSeek) assisted with data analysis, statistical computation,
visualisation, and manuscript preparation. The framework, experimental design, findings,
and interpretations remain the author's sole responsibility.

% ================================================================
\bibliographystyle{plainnat}

\begin{thebibliography}{50}

% MCH Programme papers
\bibitem[Laxman(2026a)]{laxman2026a}
Laxman, M.~M. (2026a).
\newblock Context curves behavior: Measuring {AI} relational dynamics with {$\Delta$RCI}.
\newblock \emph{Preprints.org}.
\newblock \url{https://doi.org/10.20944/preprints202601.1881.v2}

\bibitem[Laxman(2026b)]{laxman2026b}
Laxman, M.~M. (2026b).
\newblock Scaling context sensitivity: A standardized benchmark of {$\Delta$RCI} across
25 model-domain runs.
\newblock \emph{Preprints.org}.
\newblock \url{https://doi.org/10.20944/preprints202602.1114.v2}

\bibitem[Laxman(2026c)]{laxman2026c}
Laxman, M.~M. (2026c).
\newblock Domain-specific temporal dynamics of context sensitivity in large language
models.
\newblock \emph{Preprints.org}.
\newblock \url{https://doi.org/10.20944/preprints202602.1674.v1}

\bibitem[Laxman(2026d)]{laxman2026d}
Laxman, M.~M. (2026d).
\newblock Engagement as entanglement: Variance signatures of bidirectional context
coupling in large language models.
\newblock \emph{Preprints.org}.
\newblock \url{https://doi.org/10.20944/preprints202603.0055.v1}

\bibitem[Laxman(2026e)]{laxman2026e}
Laxman, M.~M. (2026e).
\newblock Stochastic incompleteness: A predictability taxonomy for clinical {AI}
deployment.
\newblock \emph{Preprints.org}.
\newblock \url{https://doi.org/10.20944/preprints202602.2034.v1}

\bibitem[Laxman(2026g)]{laxman2026g}
Laxman, M.~M. (2026g).
\newblock The structure and trajectory of context sensitivity in {LLMs}: Content-order
decomposition and variance dissociation.
\newblock \emph{Preprints.org}.
\newblock \url{https://doi.org/10.20944/preprints202603.1116.v1}

\bibitem[Laxman(2026h)]{laxman2026h}
Laxman, M.~M. (2026h).
\newblock Encoding fidelity and coherent misalignment in non-{E}nglish clinical {AI}.
\newblock \emph{Preprints.org}.
\newblock \url{https://doi.org/10.20944/preprints202604.0061.v1}

\bibitem[Laxman(2026i)]{laxman2026i}
Laxman, M.~M. (2026i).
\newblock Measurement matters: Embedding model choice determines encoding fidelity
assessment in multilingual clinical {AI}.
\newblock \emph{Zenodo}.
\newblock \url{https://doi.org/10.5281/zenodo.19466613}

% Embedding models
\bibitem[Reimers \& Gurevych(2020)]{reimers2020}
Reimers, N. \& Gurevych, I. (2020).
\newblock Making monolingual sentence embeddings multilingual using knowledge distillation.
\newblock \emph{Proceedings of EMNLP 2020}.
\newblock arXiv: 2004.09813.

\bibitem[Feng et~al.(2022)]{feng2022}
Feng, F., Yang, Y., Cer, D., Arivazhagan, N., \& Wang, W. (2022).
\newblock Language-agnostic {BERT} sentence embedding.
\newblock \emph{Proceedings of ACL 2022}.
\newblock arXiv: 2007.01852.

% Foundational
\bibitem[Shannon(1948)]{shannon1948}
Shannon, C.~E. (1948).
\newblock A mathematical theory of communication.
\newblock \emph{Bell System Technical Journal}, 27(3), 379--423.
\newblock DOI: 10.1002/j.1538-7305.1948.tb01338.x.

\bibitem[Vaswani et~al.(2017)]{vaswani2017}
Vaswani, A., Shazeer, N., Parmar, N., et~al. (2017).
\newblock Attention is all you need.
\newblock \emph{Advances in Neural Information Processing Systems}, 30.
\newblock arXiv: 1706.03762.

\bibitem[Brown et~al.(2020)]{brown2020}
Brown, T., Mann, B., Ryder, N., et~al. (2020).
\newblock Language models are few-shot learners.
\newblock \emph{Advances in Neural Information Processing Systems}, 33, 1877--1901.
\newblock arXiv: 2005.14165.

% Clinical AI
\bibitem[Singhal et~al.(2023)]{singhal2023}
Singhal, K., Azizi, S., Tu, T., et~al. (2023).
\newblock Large language models encode clinical knowledge.
\newblock \emph{Nature}, 620(7973), 172--180.
\newblock DOI: 10.1038/s41586-023-06291-2.

\bibitem[Jin et~al.(2024)]{jin2024}
Jin, Y., Chandra, M., Verma, G., Hu, Y., De~Choudhury, M., \& Kumar, S. (2024).
\newblock Better to ask in {English}: Cross-lingual evaluation of large language models
for healthcare queries.
\newblock \emph{Proceedings of the ACM Web Conference 2024}.
\newblock arXiv: 2310.13132.

\bibitem[Pfohl et~al.(2024)]{pfohl2024}
Pfohl, S.~R., Cole-Lewis, H., Sayres, R., et~al. (2024).
\newblock A toolbox for surfacing health equity harms and biases in large language models.
\newblock \emph{Nature Medicine}, 30, 3590--3600.
\newblock DOI: 10.1038/s41591-024-03258-2.

% AI safety and behaviour
\bibitem[Hubinger et~al.(2019)]{hubinger2019}
Hubinger, E., van Merwijk, C., Mikulik, V., Skalse, J., \& Garrabrant, S. (2019).
\newblock Risks from learned optimization in advanced machine learning systems.
\newblock \emph{arXiv preprint}.
\newblock arXiv: 1906.01820.

\bibitem[Shanahan(2024)]{shanahan2024}
Shanahan, M. (2024).
\newblock Talking about large language models.
\newblock \emph{Communications of the ACM}, 67(2), 68--79.
\newblock DOI: 10.1145/3624720.

\bibitem[Huang et~al.(2024)]{huang2024}
Huang, L., Yu, W., Ma, W., et~al. (2024).
\newblock A survey on hallucination in large language models: Principles, taxonomy,
challenges, and open questions.
\newblock \emph{ACM Transactions on Information Systems}, 43(2).
\newblock arXiv: 2311.05232.

% Language representation
\bibitem[Conneau et~al.(2020)]{conneau2020}
Conneau, A., Khandelwal, K., Goyal, N., et~al. (2020).
\newblock Unsupervised cross-lingual representation learning at scale.
\newblock \emph{Proceedings of ACL 2020}, 8440--8451.
\newblock arXiv: 1911.02116.

\bibitem[Blasi et~al.(2022)]{blasi2022}
Blasi, D.~E., Henrich, J., Adamou, E., Kemmerer, D., \& Majid, A. (2022).
\newblock Over-reliance on {English} hinders cognitive science.
\newblock \emph{Trends in Cognitive Sciences}, 26(12), 1153--1170.
\newblock DOI: 10.1016/j.tics.2022.09.015.

\bibitem[Kaplan et~al.(2020)]{kaplan2020}
Kaplan, J., McCandlish, S., Henighan, T., et~al. (2020).
\newblock Scaling laws for neural language models.
\newblock \emph{arXiv preprint}.
\newblock arXiv: 2001.08361.

\bibitem[Liu et~al.(2024)]{liu2024lost}
Liu, N.~F., Lin, K., Hewitt, J., Paranjape, A., Bevilacqua, M., Petroni, F.,
\& Liang, P. (2024).
\newblock Lost in the middle: How language models use long contexts.
\newblock \emph{Transactions of the Association for Computational Linguistics},
12, 157--173.
\newblock DOI: 10.1162/tacl\_a\_00638.

\bibitem[Laban et~al.(2025)]{laban2025}
Laban, P., Hayashi, H., Zhou, Y., \& Neville, J. (2025).
\newblock {LLMs} get lost in multi-turn conversation.
\newblock \emph{arXiv preprint}.
\newblock arXiv: 2505.06120.

\bibitem[Tishby \& Zaslavsky(2015)]{tishby2015}
Tishby, N. \& Zaslavsky, N. (2015).
\newblock Deep learning and the information bottleneck principle.
\newblock \emph{2015 IEEE Information Theory Workshop (ITW)}, 1--5.
\newblock DOI: 10.1109/ITW.2015.7133169.

\bibitem[Fierro et~al.(2024)]{fierro2024}
Fierro, C., Dhar, R., Stamatiou, F., Garneau, N., \& S{\o}gaard, A. (2024).
\newblock Defining knowledge: Bridging epistemology and large language models.
\newblock \emph{Proceedings of EMNLP 2024}, 16096--16111.
\newblock DOI: 10.18653/v1/2024.emnlp-main.900.

\bibitem[Kuratov et~al.(2024)]{kuratov2024}
Kuratov, Y., Bulatov, A., Anokhin, P., Rodkin, I., Sorokin, D., Sorokin, A.,
\& Burtsev, M. (2024).
\newblock {BABILong}: Testing the limits of {LLMs} with long context
reasoning-in-a-haystack.
\newblock \emph{NeurIPS 2024 Datasets and Benchmarks Track}.
\newblock arXiv: 2406.10149.

% LLM evaluation benchmarks
\bibitem[Zheng et~al.(2023)]{zheng2023}
Zheng, L., Chiang, W.-L., Sheng, Y., et~al. (2023).
\newblock Judging {LLM}-as-a-judge with {MT-Bench} and {Chatbot Arena}.
\newblock \emph{Advances in Neural Information Processing Systems 36
(NeurIPS 2023) Datasets and Benchmarks Track}.
\newblock arXiv: 2306.05685.

\bibitem[Hendrycks et~al.(2021)]{hendrycks2021}
Hendrycks, D., Burns, C., Basart, S., Zou, A., Mazeika, M., Song, D.,
\& Steinhardt, J. (2021).
\newblock Measuring massive multitask language understanding.
\newblock \emph{Proceedings of ICLR 2021}.
\newblock arXiv: 2009.03300.

% Legal AI
\bibitem[Guha et~al.(2023)]{guha2023}
Guha, N., Nyarko, J., Ho, D.~E., R{\'e}, C., Chilton, A., et~al. (2023).
\newblock {LegalBench}: A collaboratively built benchmark for measuring legal
reasoning in large language models.
\newblock \emph{Advances in Neural Information Processing Systems 36
(NeurIPS 2023) Datasets and Benchmarks Track}.
\newblock arXiv: 2308.11462.

% Sentence embeddings
\bibitem[Muennighoff et~al.(2023)]{muennighoff2023}
Muennighoff, N., Tazi, N., Magne, L., \& Reimers, N. (2023).
\newblock {MTEB}: Massive text embedding benchmark.
\newblock \emph{Proceedings of EACL 2023}, 2014--2037.
\newblock DOI: 10.18653/v1/2023.eacl-main.148.

\bibitem[Gao et~al.(2021)]{gao2021simcse}
Gao, T., Yao, X., \& Chen, D. (2021).
\newblock {SimCSE}: Simple contrastive learning of sentence embeddings.
\newblock \emph{Proceedings of EMNLP 2021}, 6894--6910.
\newblock DOI: 10.18653/v1/2021.emnlp-main.552.

% Alignment
\bibitem[Ouyang et~al.(2022)]{ouyang2022}
Ouyang, L., Wu, J., Jiang, X., Almeida, D., Wainwright, C.~L.,
Mishkin, P., et~al. (2022).
\newblock Training language models to follow instructions with human feedback.
\newblock \emph{Advances in Neural Information Processing Systems 35
(NeurIPS 2022)}.
\newblock arXiv: 2203.02155.

% Medical AI
\bibitem[Thirunavukarasu et~al.(2023)]{thirunavukarasu2023}
Thirunavukarasu, A.~J., Ting, D.~S.~J., Elangovan, K., et~al. (2023).
\newblock Large language models in medicine.
\newblock \emph{Nature Medicine}, 29(8), 1930--1940.
\newblock DOI: 10.1038/s41591-023-02448-8.

\bibitem[Nori et~al.(2023)]{nori2023}
Nori, H., Lee, Y.~T., Zhang, S., et~al. (2023).
\newblock Can generalist foundation models outcompete special-purpose tuning?
{C}ase study in medicine.
\newblock \emph{arXiv preprint}.
\newblock arXiv: 2311.16452.

% Multilingual evaluation
\bibitem[Ahuja et~al.(2023)]{ahuja2023}
Ahuja, K., Diddee, H., Hada, R., et~al. (2023).
\newblock {MEGA}: Multilingual evaluation of generative {AI}.
\newblock \emph{Proceedings of EMNLP 2023}.
\newblock DOI: 10.18653/v1/2023.emnlp-main.258.

% Scaling and architecture
\bibitem[Hoffmann et~al.(2022)]{hoffmann2022}
Hoffmann, J., Borgeaud, S., Mensch, A., et~al. (2022).
\newblock Training compute-optimal large language models.
\newblock \emph{Advances in Neural Information Processing Systems 35
(NeurIPS 2022)}.
\newblock arXiv: 2203.15556.

% Output variance and reproducibility
\bibitem[Atil et~al.(2024)]{atil2024}
Atil, B., Aykent, S., Chittams, A., et~al. (2024).
\newblock Non-determinism of ``deterministic'' {LLM} settings.
\newblock \emph{arXiv preprint}.
\newblock arXiv: 2408.04667.

\bibitem[Renze \& Guven(2024)]{renze2024}
Renze, M. \& Guven, E. (2024).
\newblock The effect of sampling temperature on problem solving in large
language models.
\newblock \emph{Findings of EMNLP 2024}.
\newblock DOI: 10.18653/v1/2024.findings-emnlp.432.

% Information bottleneck
\bibitem[Saxe et~al.(2019)]{saxe2019}
Saxe, A., Bansal, Y., Dapello, J., et~al. (2019).
\newblock On the information bottleneck theory of deep learning.
\newblock \emph{Journal of Statistical Mechanics: Theory and Experiment},
2019(12), 124020.
\newblock DOI: 10.1088/1742-5468/ab3985.

\bibitem[Shwartzziv \& Tishby(2017)]{shwartzziv2017}
Shwartz-Ziv, R. \& Tishby, N. (2017).
\newblock Opening the black box of deep neural networks via information.
\newblock \emph{arXiv preprint}.
\newblock arXiv: 1703.00810.

% AI safety and hallucination
\bibitem[Ji et~al.(2023)]{ji2023}
Ji, Z., Lee, N., Frieske, R., et~al. (2023).
\newblock Survey of hallucination in natural language generation.
\newblock \emph{ACM Computing Surveys}, 55(12), 1--38.
\newblock DOI: 10.1145/3571730.

\bibitem[Lin et~al.(2022)]{lin2022}
Lin, S., Hilton, J., \& Evans, O. (2022).
\newblock {TruthfulQA}: Measuring how models mimic human falsehoods.
\newblock \emph{Proceedings of ACL 2022}, 3214--3252.
\newblock DOI: 10.18653/v1/2022.acl-long.229.

% Epistemological approaches
\bibitem[Bender et~al.(2021)]{bender2021}
Bender, E.~M., Gebru, T., McMillan-Major, A., \& Shmitchell, S. (2021).
\newblock On the dangers of stochastic parrots: {C}an language models be
too big?
\newblock \emph{Proceedings of FAccT 2021}, 610--623.
\newblock DOI: 10.1145/3442188.3445922.

\bibitem[Chang et~al.(2024)]{chang2024survey}
Chang, Y., Wang, X., Wang, J., et~al. (2024).
\newblock A survey on evaluation of large language models.
\newblock \emph{ACM Transactions on Intelligent Systems and Technology},
15(3), 1--45.
\newblock DOI: 10.1145/3641289.

% Multi-turn and context evaluation
\bibitem[Wang et~al.(2024a)]{wang2024mint}
Wang, X., Wang, Z., Liu, J., Chen, Y., Yuan, L., Peng, H., \& Ji, H. (2024).
\newblock {MINT}: Evaluating {LLMs} in multi-turn interaction with tools and language feedback.
\newblock \emph{ICLR 2024}.
\newblock arXiv: 2309.10691.

\bibitem[Kwan et~al.(2024)]{kwan2024mteval}
Kwan, W.-C., Zeng, X., Jiang, Y., Wang, Y., Li, L., Shang, L., Jiang, X., Liu, Q., \& Wong, K.-F. (2024).
\newblock {MT-Eval}: A multi-turn capabilities evaluation benchmark for large language models.
\newblock \emph{Proceedings of EMNLP 2024}, 20153--20177.
\newblock DOI: 10.18653/v1/2024.emnlp-main.1124.

\bibitem[Katsis et~al.(2025)]{katsis2025mtrag}
Katsis, Y., Rosenthal, S., Fadnis, K., et~al. (2025).
\newblock {mtRAG}: A multi-turn conversational benchmark for evaluating retrieval-augmented generation systems.
\newblock \emph{Transactions of the Association for Computational Linguistics}, 13, 784--808.
\newblock DOI: 10.1162/TACL.a.19.

\bibitem[Bai et~al.(2024a)]{bai2024longbench}
Bai, Y., Lv, X., Zhang, J., et~al. (2024).
\newblock {LongBench}: A bilingual, multitask benchmark for long context understanding.
\newblock \emph{Proceedings of ACL 2024}, 3119--3137.
\newblock arXiv: 2308.14508.

\bibitem[Zhang et~al.(2024a)]{zhang2024infinitebench}
Zhang, X., Chen, Y., Hu, S., et~al. (2024).
\newblock {InfiniteBench}: Extending long context evaluation beyond 100{K} tokens.
\newblock \emph{Proceedings of ACL 2024}, 15262--15277.
\newblock arXiv: 2402.13718.

\bibitem[Hsieh et~al.(2024)]{hsieh2024ruler}
Hsieh, C.-P., Sun, S., Kriman, S., et~al. (2024).
\newblock {RULER}: What's the real context size of your long-context language models?
\newblock \emph{arXiv preprint}.
\newblock arXiv: 2404.06654.

\bibitem[An et~al.(2025)]{an2025context}
An, C., Zhang, J., Zhong, M., et~al. (2025).
\newblock Why does the effective context length of {LLMs} fall short?
\newblock \emph{ICLR 2025}.
\newblock arXiv: 2410.18745.

\bibitem[Maharana et~al.(2024)]{maharana2024memory}
Maharana, A., Lee, D.-H., Tulyakov, S., Bansal, M., Barbieri, F., \& Fang, Y. (2024).
\newblock Evaluating very long-term conversational memory of {LLM} agents.
\newblock \emph{Proceedings of ACL 2024}, 13851--13870.
\newblock arXiv: 2402.17753.

% Systematic LLM evaluation
\bibitem[Liang et~al.(2023)]{liang2023helm}
Liang, P., Bommasani, R., Lee, T., et~al. (2023).
\newblock Holistic evaluation of language models.
\newblock \emph{Transactions on Machine Learning Research}.
\newblock arXiv: 2211.09110.

\bibitem[Chiang et~al.(2024)]{chiang2024chatbot}
Chiang, W.-L., Zheng, L., Sheng, Y., et~al. (2024).
\newblock Chatbot {Arena}: An open platform for evaluating {LLMs} by human preference.
\newblock \emph{Proceedings of ICML 2024}, vol.~235, 8359--8388.
\newblock arXiv: 2403.04132.

\bibitem[Sclar et~al.(2024)]{sclar2024sensitivity}
Sclar, M., Choi, Y., Tsvetkov, Y., \& Suhr, A. (2024).
\newblock Quantifying language models' sensitivity to spurious features in prompt design.
\newblock \emph{ICLR 2024}.
\newblock arXiv: 2310.11324.

% Medical AI
\bibitem[Singhal et~al.(2025)]{singhal2025expert}
Singhal, K., Tu, T., Gottweis, J., et~al. (2025).
\newblock Toward expert-level medical question answering with large language models.
\newblock \emph{Nature Medicine}, 31, 943--950.
\newblock DOI: 10.1038/s41591-024-03423-7.

\bibitem[Ayers et~al.(2023)]{ayers2023physician}
Ayers, J.~W., Poliak, A., Dredze, M., et~al. (2023).
\newblock Comparing physician and artificial intelligence chatbot responses to patient questions.
\newblock \emph{JAMA Internal Medicine}, 183(6), 589--596.
\newblock DOI: 10.1001/jamainternmed.2023.1838.

% Legal AI
\bibitem[Nay et~al.(2024)]{nay2024tax}
Nay, J.~J., Karamardian, D., Lawsky, S.~B., et~al. (2024).
\newblock Large language models as tax attorneys: A case study in legal capabilities emergence.
\newblock \emph{Philosophical Transactions of the Royal Society A}, 382(2270).
\newblock DOI: 10.1098/rsta.2023.0159.

\bibitem[Dahl et~al.(2024)]{dahl2024hallucinations}
Dahl, M., Magesh, V., Suzgun, M., \& Ho, D.~E. (2024).
\newblock Large legal fictions: Profiling legal hallucinations in large language models.
\newblock \emph{Journal of Legal Analysis}, 16(1), 64--93.
\newblock DOI: 10.1093/jla/laae003.

\bibitem[Choi et~al.(2022)]{choi2022lawschool}
Choi, J.~H., Hickman, K.~E., Monahan, A., \& Schwarcz, D. (2022).
\newblock {ChatGPT} goes to law school.
\newblock \emph{Journal of Legal Education}, 71(3).
\newblock SSRN: 4335905.

% Epistemological and Ethics
\bibitem[Suzgun et~al.(2024)]{suzgun2024belief}
Suzgun, M., Gur, T., Bianchi, F., et~al. (2024).
\newblock Belief in the machine: Investigating epistemological blind spots of language models.
\newblock \emph{arXiv preprint}.
\newblock arXiv: 2410.21195.

\bibitem[Paseri \& Durante(2025)]{paseri2025epistemological}
Paseri, L. \& Durante, M. (2025).
\newblock Examining epistemological challenges of large language models in law.
\newblock \emph{Cambridge Forum on AI: Law and Governance}, 1, e7.
\newblock DOI: 10.1017/aig.2024.7.

\bibitem[Ji et~al.(2024)]{ji2024moralbench}
Ji, J., Chen, Y., Jin, M., et~al. (2024).
\newblock {MoralBench}: Moral evaluation of {LLMs}.
\newblock \emph{arXiv preprint}.
\newblock arXiv: 2406.04428.

% Output variance and stochasticity
\bibitem[Shi et~al.(2024)]{shi2024decoding}
Shi, C., Yang, H., Cai, D., et~al. (2024).
\newblock A thorough examination of decoding methods in the era of {LLMs}.
\newblock \emph{Proceedings of EMNLP 2024}, 8601--8629.
\newblock DOI: 10.18653/v1/2024.emnlp-main.489.

\bibitem[Dietrich \& Hollstein(2025)]{dietrich2025reproducibility}
Dietrich, J. \& Hollstein, A. (2025).
\newblock Performance and reproducibility of large language models in named entity recognition.
\newblock \emph{Drug Safety}, 48(3), 287--303.
\newblock DOI: 10.1007/s40264-024-01499-1.

% Sentence embeddings
\bibitem[Steck et~al.(2024)]{steck2024cosine}
Steck, H., Ekanadham, C., \& Kallus, N. (2024).
\newblock Is cosine-similarity of embeddings really about similarity?
\newblock \emph{Companion Proceedings of the ACM Web Conference 2024}.
\newblock DOI: 10.1145/3589335.3651526.

\bibitem[Ethayarajh(2019)]{ethayarajh2019contextual}
Ethayarajh, K. (2019).
\newblock How contextual are contextualized word representations?
\newblock \emph{Proceedings of EMNLP-IJCNLP 2019}, 55--65.
\newblock DOI: 10.18653/v1/D19-1006.

\bibitem[Chen et~al.(2024a)]{chen2024m3embedding}
Chen, J., Xiao, S., Zhang, P., Luo, K., Lian, D., \& Liu, Z. (2024).
\newblock {M3-Embedding}: Multi-linguality, multi-functionality, multi-granularity text embeddings through self-knowledge distillation.
\newblock \emph{Findings of ACL 2024}.
\newblock DOI: 10.18653/v1/2024.findings-acl.137.

\bibitem[Enevoldsen et~al.(2025)]{enevoldsen2025mmteb}
Enevoldsen, K., Chung, I., Kerboua, I., et~al. (2025).
\newblock {MMTEB}: Massive multilingual text embedding benchmark.
\newblock \emph{ICLR 2025}.
\newblock arXiv: 2502.13595.

% Information theory and scaling
\bibitem[Kawaguchi et~al.(2023)]{kawaguchi2023bottleneck}
Kawaguchi, K., Deng, Z., Ji, X., \& Huang, J. (2023).
\newblock How does information bottleneck help deep learning?
\newblock \emph{Proceedings of ICML 2023}, 16049--16096.
\newblock arXiv: 2305.18887.

\bibitem[Bahri et~al.(2024)]{bahri2024scaling}
Bahri, Y., Dyer, E., Kaplan, J., Lee, J., \& Sharma, U. (2024).
\newblock Explaining neural scaling laws.
\newblock \emph{Proceedings of the National Academy of Sciences}, 121(27).
\newblock DOI: 10.1073/pnas.2311878121.

% AI safety and trustworthiness
\bibitem[Wang et~al.(2023)]{wang2023decodingtrust}
Wang, B., Chen, W., Pei, H., et~al. (2023).
\newblock {DecodingTrust}: A comprehensive assessment of trustworthiness in {GPT} models.
\newblock \emph{NeurIPS 2023}.
\newblock arXiv: 2306.11698.

\bibitem[Sharma et~al.(2024)]{sharma2024sycophancy}
Sharma, M., Tong, M., Korbak, T., et~al. (2024).
\newblock Towards understanding sycophancy in language models.
\newblock \emph{ICLR 2024}.
\newblock arXiv: 2310.13548.

\bibitem[Manakul et~al.(2023)]{manakul2023selfcheck}
Manakul, P., Liusie, A., \& Gales, M. (2023).
\newblock {SelfCheckGPT}: Zero-resource black-box hallucination detection for generative large language models.
\newblock \emph{Proceedings of EMNLP 2023}, 9004--9017.
\newblock DOI: 10.18653/v1/2023.emnlp-main.557.

% Transformer architecture
\bibitem[Fedus et~al.(2022)]{fedus2022switch}
Fedus, W., Zoph, B., \& Shazeer, N. (2022).
\newblock Switch transformers: Scaling to trillion parameter models with simple and efficient sparsity.
\newblock \emph{Journal of Machine Learning Research}, 23(120), 1--40.

\bibitem[Wei et~al.(2022a)]{wei2022cot}
Wei, J., Wang, X., Schuurmans, D., et~al. (2022).
\newblock Chain-of-thought prompting elicits reasoning in large language models.
\newblock \emph{NeurIPS 2022}.
\newblock arXiv: 2201.11903.

\bibitem[Olsson et~al.(2022)]{olsson2022induction}
Olsson, C., Elhage, N., Nanda, N., et~al. (2022).
\newblock In-context learning and induction heads.
\newblock \emph{Transformer Circuits Thread}.
\newblock arXiv: 2209.11895.

\bibitem[Wei et~al.(2022b)]{wei2022emergent}
Wei, J., Tay, Y., Bommasani, R., et~al. (2022).
\newblock Emergent abilities of large language models.
\newblock \emph{Transactions on Machine Learning Research}.
\newblock arXiv: 2206.07682.

\bibitem[Bai et~al.(2022)]{bai2022constitutional}
Bai, Y., Kadavath, S., Kundu, S., et~al. (2022).
\newblock Constitutional {AI}: Harmlessness from {AI} feedback.
\newblock \emph{arXiv preprint}.
\newblock arXiv: 2212.08073.

% Multilingual
\bibitem[Lai et~al.(2023)]{lai2023chatgpt}
Lai, V.~D., Ngo, N., Veyseh, A.~P.~B., et~al. (2023).
\newblock {ChatGPT} beyond {English}: Towards a comprehensive evaluation of large language models in multilingual learning.
\newblock \emph{Findings of EMNLP 2023}, 13171--13189.
\newblock DOI: 10.18653/v1/2023.findings-emnlp.878.

\bibitem[Joshi et~al.(2020)]{joshi2020diversity}
Joshi, P., Santy, S., Budhiraja, A., Bali, K., \& Choudhury, M. (2020).
\newblock The state and fate of linguistic diversity and inclusion in the {NLP} world.
\newblock \emph{Proceedings of ACL 2020}.
\newblock DOI: 10.18653/v1/2020.acl-main.560.

\bibitem[Gala et~al.(2023)]{gala2023indictrans2}
Gala, J., Chitale, P.~A., et~al. (2023).
\newblock {IndicTrans2}: Towards high-quality and accessible machine translation models for all 22 scheduled {Indian} languages.
\newblock \emph{Transactions on Machine Learning Research}.
\newblock arXiv: 2305.16307.

\end{thebibliography}

\end{document}
